% =============================================================================
%  Jacobian Derivation for a Hybrid 2-DoF Planar Cable-Driven Manipulator
%  -----------------------------------------------------------------------------
%  System: CupManipulatorTendon
%    - Joint 1 (link1_base,   q1): directly actuated revolute joint
%    - Joint 2 (link2_link1,  q2): cable-driven via HTD 5M 60T belt/pulley
%  Robot moves in the horizontal XY plane (world Z = constant).
%  Link lengths (from FK at q = 0):
%    L1 = 342.47 mm   (joint-1 origin → joint-2 origin)
%    L2 = 190.00 mm   (joint-2 origin → end-effector)
%  Pulley parameters (HTD 5M 60T):
%    Pitch = 5 mm,  Teeth = 60  =>  r_p = 60 x 5 / (2 pi) ≈ 47.75 mm
% =============================================================================

\documentclass[12pt, a4paper]{article}

% ── Packages ─────────────────────────────────────────────────────────────────
\usepackage{amsmath, amssymb, bm}
\usepackage{geometry}
\usepackage{graphicx}
\usepackage{array, booktabs}
\usepackage{hyperref}
\usepackage{xcolor}
\usepackage{mathtools}
\usepackage{cleveref}
\usepackage{enumitem}
\usepackage{amsthm}
\usepackage[backgroundcolor=blue!5, linecolor=blue!40, linewidth=1pt,
            innertopmargin=7pt, innerbottommargin=7pt,
            innerleftmargin=8pt, innerrightmargin=8pt]{mdframed}

\theoremstyle{definition}
\newtheorem{corollary}{Corollary}[section]

\geometry{margin=2.5cm}

% ── Macros ───────────────────────────────────────────────────────────────────
% Bold roman vectors / matrices
\newcommand{\bq}{\bm{q}}
\newcommand{\bp}{\bm{p}}
\newcommand{\bu}{\bm{u}}
\newcommand{\btau}{\bm{\tau}}
\newcommand{\bF}{\bm{F}}
\newcommand{\bJ}{\bm{J}}
\newcommand{\bA}{\bm{A}}
\newcommand{\bM}{\bm{M}}
\newcommand{\bC}{\bm{C}}
\newcommand{\bG}{\bm{G}}
\newcommand{\bl}{\bm{l}}
\newcommand{\bomega}{\bm{\omega}}

% Short-hand trig
\newcommand{\cq}[1]{c_{#1}}   % cos(q_#1)
\newcommand{\sq}[1]{s_{#1}}   % sin(q_#1)
\newcommand{\cqq}[2]{c_{{#1}{#2}}}  % cos(q_#1 + q_#2)
\newcommand{\sqq}[2]{s_{{#1}{#2}}}  % sin(q_#1 + q_#2)

% Frames
\newcommand{\FO}{\{F_O\}}      % inertial / world frame
\newcommand{\FE}{\{F_\varepsilon\}} % EE frame

% =============================================================================
\begin{document}

\title{%
  Jacobian Derivation for a Hybrid 2-DoF\\
  Planar Cable-Driven Manipulator%
}
\author{CupManipulatorTendon --- Kinematic \& Dynamic Analysis}
\date{March 2026}
\maketitle
\tableofcontents
\newpage

% =============================================================================
\section{System Description}
\label{sec:system}
% =============================================================================

The manipulator under consideration is a \textbf{hybrid 2-DoF planar serial arm}
whose end-effector (EE) moves in a fixed horizontal plane (world $Z = z_0$).  The
two degrees of freedom are:

\begin{itemize}[leftmargin=2em]
  \item $q_1$ (\texttt{link1\_base}): revolute joint, \emph{directly actuated} by
        motor torque $\tau_1$.
  \item $q_2$ (\texttt{link2\_link1}): revolute joint, \emph{cable-driven} via a
        HTD 5M belt wound on a 60-tooth pulley.  The controlled input is the cable
        tension $f_c \ge 0$.
\end{itemize}

\noindent
This hybrid actuation is analogous to the Cable-Driven Parallel Manipulator (CDPM)
framework presented in \cite{pott2018cable}, where a structure matrix maps cable
forces to generalised joint torques.  Here, only joint~2 is governed by cable
mechanics; joint~1 retains the classical direct-drive model.

\paragraph{Geometric parameters.}
\begin{center}
  \renewcommand{\arraystretch}{1.3}
  \begin{tabular}{lll}
    \toprule
    Symbol & Value & Description \\
    \midrule
    $L_1$ & 342.47 mm & Distance from joint-1 axis to joint-2 axis (horizontal) \\
    $L_2$ & 190.00 mm & Distance from joint-2 axis to EE origin (horizontal) \\
    $r_p$ & 47.75 mm  & HTD 5M 60T pulley pitch radius ($= 60 \times 5\,\text{mm} / (2\pi)$) \\
    $z_0$ & 321.55 mm & Constant world-Z of the EE (set at home pose) \\
    \bottomrule
  \end{tabular}
\end{center}

% =============================================================================
\section{Coordinate Frames and Notation}
\label{sec:frames}
% =============================================================================

Following the CDPM convention in the attached reference:

\begin{itemize}
  \item $\FO$ --- inertial (world) frame; origin at joint-1 axis.
  \item $\FE$ --- end-effector body frame; rigidly attached to link~2, origin at
        the EE point.
  \item Cable attachment on the \emph{base} (drive pulley centre):
        ${}^O\bm{r}_{OA}$ --- fixed in $\FO$.
  \item Cable attachment on the \emph{EE link} (driven pulley winding point):
        ${}^\varepsilon\bm{r}_{OB}$ --- fixed in $\FE$.
\end{itemize}

\noindent
Notation conventions:
$c_i \triangleq \cos q_i$,\;
$s_i \triangleq \sin q_i$,\;
$c_{12} \triangleq \cos(q_1 + q_2)$,\;
$s_{12} \triangleq \sin(q_1 + q_2)$.

% =============================================================================
\section{Forward Kinematics}
\label{sec:fk}
% =============================================================================

The EE position in world frame $\FO$ is:

\begin{equation}
  \label{eq:fk}
  \bp(q_1, q_2) =
  \begin{bmatrix}
    p_x \\ p_y
  \end{bmatrix}
  =
  \begin{bmatrix}
    L_1 \cq{1} + L_2 \cqq{1}{2} \\
    L_1 \sq{1} + L_2 \sqq{1}{2}
  \end{bmatrix},
\end{equation}

where only the horizontal components are shown; $p_z = z_0 = \text{const}$.

% =============================================================================
\section{Geometric (Translational) Jacobian}
\label{sec:geo_jacobian}
% =============================================================================

Differentiating \cref{eq:fk} with respect to $\bq = [q_1, q_2]^\top$ gives the
standard geometric Jacobian $\bJ \in \mathbb{R}^{2 \times 2}$:

\begin{equation}
  \label{eq:Jgeo}
  \bJ(\bq)
  = \frac{\partial \bp}{\partial \bq}
  = \begin{bmatrix}
      -L_1 \sq{1} - L_2 \sqq{1}{2} & -L_2 \sqq{1}{2} \\[4pt]
       L_1 \cq{1} + L_2 \cqq{1}{2} &  L_2 \cqq{1}{2}
    \end{bmatrix}.
\end{equation}

\noindent
Velocity-level kinematics:
\begin{equation}
  \label{eq:vel_fk}
  \dot{\bp} = \bJ(\bq)\,\dot{\bq}.
\end{equation}

\paragraph{Singularity analysis.}
The determinant of $\bJ$ is:
\begin{equation}
  \det\bJ = L_1 L_2 \sin q_2 .
\end{equation}
A kinematic singularity occurs whenever $q_2 = 0$ or $q_2 = \pm\pi$
(fully-extended or folded configurations).  This is the fundamental reason that
Jacobian-based (velocity) IK is unreliable near $q_2 = 0$ and analytical 2R IK
must be used as a fallback.

% =============================================================================
\section{Actuation Models}
\label{sec:actuation}
% =============================================================================

\subsection{Joint 1 --- Direct Actuation}
\label{sec:act_j1}

Joint~1 is driven directly: the motor output torque $\tau_1$ appears
unmodified as the generalised force for $q_1$.  There is no cable loop;
the actuation ratio is unity.

\subsection{Joint 2 --- Cable (Belt/Pulley) Actuation}
\label{sec:act_j2}

Following the CDPM cable-vector formulation \cite{pott2018cable}, a single cable
wraps around a drive pulley (radius $r_p$) mounted on the base and is anchored to
a driven pulley fixed to link~2.  The cable vector ${}^\varepsilon\bm{l}$ expressed
in frame $\FE$ is (cf.\ equation~3.7 in the reference):

\begin{equation}
  \label{eq:cable_vec}
  {}^\varepsilon\bm{l}
  = -\,{}_O^\varepsilon R\,{}^O\bm{r}_{OA} + {}^\varepsilon\bm{r}_{OB},
\end{equation}

where ${}_O^\varepsilon R$ is the rotation from $\FO$ to $\FE$ and ${}^O\bm{r}_{OA}$
is the fixed world-frame position of the cable's base-side anchor.  The cable
length is $l = \|{}^\varepsilon\bm{l}\|$.

\paragraph{Belt--pulley simplification.}
Because both the drive and driven pulleys are of the same HTD-5M tooth profile and
the belt is inextensible, the cable wraps without slip.  The arc-length constraint
gives a \emph{linear} relationship between cable displacement $\Delta l$ and joint
angle increment $\Delta q_2$:

\begin{equation}
  \label{eq:cable_kin}
  \dot{l} = r_p\,\dot{q}_2 ,
  \qquad
  \Delta l = r_p\,\Delta q_2 .
\end{equation}

This replaces the full CDPM cable-length formula: the wrapped geometry reduces to a
scalar transmission ratio $r_p$.

\paragraph{Cable tension and joint torque.}
The cable tension $f_c \ge 0$ produces a torque about joint~2's axis via the moment
arm $r_p$:
\begin{equation}
  \label{eq:cable_torque}
  \tau_2 = r_p\,f_c .
\end{equation}

% =============================================================================
\section{Actuation Mapping Matrix}
\label{sec:act_map}
% =============================================================================

Collect the actuator inputs into the vector
$\bu = [\tau_1,\; f_c]^\top \in \mathbb{R}^2$ and the joint torques into
$\btau = [\tau_1,\; \tau_2]^\top$.  From \cref{eq:cable_torque}:

\begin{equation}
  \label{eq:act_map}
  \btau = \bA\,\bu,
  \qquad
  \bA =
  \begin{bmatrix}
    1 & 0 \\
    0 & r_p
  \end{bmatrix},
\end{equation}

where $\bA \in \mathbb{R}^{2 \times 2}$ is the \emph{actuation matrix}.
Because $\bA$ is diagonal and invertible (for $r_p > 0$), the system is
fully actuated.

The inverse --- mapping joint torques to actuator inputs --- is:
\begin{equation}
  \bA^{-1} =
  \begin{bmatrix}
    1 & 0 \\
    0 & 1/r_p
  \end{bmatrix}.
\end{equation}

% =============================================================================
\section{Hybrid Jacobian Derivation}
\label{sec:hybrid_jac}
% =============================================================================

\subsection{Task-Space Jacobian (Velocity Level)}
\label{sec:task_vel}

From \cref{eq:vel_fk,eq:cable_kin}, the joint velocities $\dot{\bq}$ relate to
both the Cartesian EE velocity $\dot{\bp}$ and the cable rate $\dot{l}$:

\begin{equation}
  \label{eq:augmented}
  \begin{bmatrix}
    \dot{\bp} \\ \dot{l}
  \end{bmatrix}
  =
  \underbrace{%
    \begin{bmatrix}
      \bJ(\bq) \\[2pt]
      \bm{0}^\top & r_p
    \end{bmatrix}
  }_{J_{\text{aug}} \;\in\; \mathbb{R}^{3 \times 2}}
  \dot{\bq},
\end{equation}

where $\bm{0}^\top = [0,\,0]$ pads $\bJ$ to accommodate the cable row.

\noindent
Explicitly:
\begin{equation}
  \label{eq:Jaug_explicit}
  J_{\text{aug}} =
  \begin{bmatrix}
    -L_1 \sq{1} - L_2 \sqq{1}{2} & -L_2 \sqq{1}{2} \\[4pt]
     L_1 \cq{1} + L_2 \cqq{1}{2} &  L_2 \cqq{1}{2} \\[4pt]
     0 & r_p
  \end{bmatrix}.
\end{equation}

\subsection{Actuator-Space Jacobian}
\label{sec:act_jac}

Define the \emph{actuator velocity vector}
$\dot{\bu}_{v} = [\dot{q}_1,\; \dot{l}]^\top$ (motor angular rate and cable
linear rate).  The joint velocities are recovered via $\bA^{-1}$:

\begin{equation}
  \label{eq:q_from_uv}
  \dot{\bq}
  = \bA^{-1}\,\dot{\bu}_v
  =
  \begin{bmatrix}
    1 & 0 \\ 0 & 1/r_p
  \end{bmatrix}
  \begin{bmatrix}
    \dot{q}_1 \\ \dot{l}
  \end{bmatrix}.
\end{equation}

Substituting into \cref{eq:vel_fk}:
\begin{equation}
  \label{eq:hybrid_vel}
  \dot{\bp}
  = \bJ(\bq)\,\bA^{-1}\,\dot{\bu}_v
  = \underbrace{\bJ(\bq)\,\bA^{-1}}_{\bJ_h(\bq)}\,\dot{\bu}_v,
\end{equation}

\begin{equation}
  \label{eq:Jh}
  \boxed{
  \bJ_h(\bq)
  \;=\;
  \bJ(\bq)\,\bA^{-1}
  \;=\;
  \begin{bmatrix}
    -(L_1 \sq{1} + L_2 \sqq{1}{2}) & -\dfrac{L_2 \sqq{1}{2}}{r_p} \\[8pt]
     L_1 \cq{1} + L_2 \cqq{1}{2}  &  \dfrac{L_2 \cqq{1}{2}}{r_p}
  \end{bmatrix}
  }
\end{equation}

$\bJ_h \in \mathbb{R}^{2 \times 2}$ is the \textbf{hybrid Jacobian}: it maps the
\emph{physical actuator rates} (joint-1 angular velocity, cable linear velocity)
to the Cartesian EE velocity.

\paragraph{Interpretation.}
The first column of $\bJ_h$ is identical to the first column of $\bJ$: rotating
joint~1 at rate $\dot{q}_1$ sweeps the EE through the same arc regardless of the
drive type.  The second column is the standard cable column scaled by $1/r_p$,
converting from cable linear speed to the equivalent joint angular speed.

\subsection{Singularity of the Hybrid Jacobian}
\label{sec:sing}

\begin{equation}
  \det \bJ_h = \frac{1}{r_p}\det\bJ = \frac{L_1 L_2}{r_p}\sin q_2 .
\end{equation}

Since $r_p > 0$, $\bJ_h$ shares the same singularity locus as the geometric
Jacobian: $q_2 = 0\pmod{\pi}$.

% =============================================================================
\section{Cable-Force Jacobian and Statics}
\label{sec:cable_force_jac}
% =============================================================================

Following the CDPM structure-matrix formulation (equation~3.15 in the reference),
the transpose of the Jacobian maps Cartesian EE forces to actuator inputs via
virtual work:

\begin{equation}
  \delta W
  = \bF_{EE}^\top \delta\bp
  = \btau^\top \delta\bq
  = \bu^\top \bA^\top \delta\bq .
\end{equation}

Since $\delta\bp = \bJ\,\delta\bq$:
\begin{equation}
  \bF_{EE}^\top \bJ\,\delta\bq
  = \bu^\top \bA^\top \delta\bq
  \;\implies\;
  \bJ^\top \bF_{EE} = \bA^\top \bu,
\end{equation}

\begin{equation}
  \label{eq:force_map}
  \bu = (\bA^\top)^{-1} \bJ^\top \bF_{EE}
      = \bJ_h^\top \bF_{EE}.
\end{equation}

\noindent
Explicitly:
\begin{equation}
  \label{eq:Jh_T}
  \bJ_h^\top =
  \begin{bmatrix}
    -(L_1 \sq{1} + L_2 \sqq{1}{2}) & L_1 \cq{1} + L_2 \cqq{1}{2} \\[6pt]
    -\dfrac{L_2 \sqq{1}{2}}{r_p}   & \dfrac{L_2 \cqq{1}{2}}{r_p}
  \end{bmatrix}.
\end{equation}

The first row gives the required motor torque $\tau_1$; the second row gives the
required cable tension $f_c$ to produce the desired EE wrench $\bF_{EE}$.

% =============================================================================
\section{Velocity-Level Inverse Kinematics}
\label{sec:vel_ik}
% =============================================================================

\subsection{Standard Damped Pseudo-Inverse on \texorpdfstring{$\bJ$}{J}}

Given a desired EE velocity $\dot{\bp}_{\text{des}}$, the joint velocity is:
\begin{equation}
  \label{eq:q_dot_ik}
  \dot{\bq}
  = \bJ^\dag\,\dot{\bp}_{\text{des}}
  = \bJ^\top\!\left(\bJ \bJ^\top + \lambda \bm{I}_2\right)^{-1}\dot{\bp}_{\text{des}},
\end{equation}

where $\lambda > 0$ is the damping factor (Levenberg--Marquardt regularisation).
This is the method implemented in \texttt{compute\_velocity\_ik}.

\subsection{Hybrid-Jacobian Velocity IK}

To work directly in actuator space (i.e., command motor velocity $\dot{q}_1$ and
cable rate $\dot{l}$):

\begin{equation}
  \label{eq:uv_ik}
  \dot{\bu}_v
  = \bJ_h^\dag\,\dot{\bp}_{\text{des}}
  = \bJ_h^\top\!\left(\bJ_h \bJ_h^\top + \lambda \bm{I}_2\right)^{-1}
    \dot{\bp}_{\text{des}}.
\end{equation}

Then the cable command is $\dot{l}$ and the motor-1 command is $\dot{q}_1 =
[\dot{\bu}_v]_1$.

\subsection{Equivalence}

Because $\bA^{-1}$ is invertible:
\begin{equation}
  \bJ_h^\dag = \bA^{-1} \bJ^\dag \quad (\text{when }\bA \text{ is invertible}).
\end{equation}
The two velocity-IK expressions \cref{eq:q_dot_ik,eq:uv_ik} yield the same EE
motion; they differ only in which actuator space the solution is expressed.

% =============================================================================
\section{Equations of Motion}
\label{sec:eom}
% =============================================================================

The standard manipulator dynamics in joint space:
\begin{equation}
  \label{eq:eom}
  \bM(\bq)\ddot{\bq} + \bC(\bq, \dot{\bq})\dot{\bq} + \bG(\bq)
  = \btau + \bJ^\top \bF_{\text{ext}},
\end{equation}

where $\bM$ is the mass-inertia matrix, $\bC$ the Coriolis/centrifugal matrix, and
$\bG$ the gravitational vector.  Substituting $\btau = \bA\,\bu$ from
\cref{eq:act_map}:

\begin{equation}
  \boxed{
  \bM(\bq)\ddot{\bq} + \bC(\bq, \dot{\bq})\dot{\bq} + \bG(\bq)
  = \bA\,\bu + \bJ^\top \bF_{\text{ext}} .
  }
\end{equation}

For torque control, the desired actuator inputs are computed as:
\begin{equation}
  \bu = \bA^{-1}
        \bigl(\bM(\bq)\ddot{\bq}_{\text{ref}}
              + \bC(\bq,\dot{\bq})\dot{\bq}_{\text{ref}}
              + \bG(\bq)
              - \bJ^\top \bF_{\text{ext}}\bigr).
\end{equation}

% =============================================================================
\section{Summary}
\label{sec:summary}
% =============================================================================

\begin{center}
  \renewcommand{\arraystretch}{1.5}
  \begin{tabular}{lll}
    \toprule
    Symbol & Expression & Role \\
    \midrule
    $\bJ(\bq)$       & \cref{eq:Jgeo} & Geometric Jacobian ($2\times 2$)\\
    $\bA$            & \cref{eq:act_map} & Actuation matrix (diagonal, $2\times 2$)\\
    $\bJ_h = \bJ\bA^{-1}$ & \cref{eq:Jh} & Hybrid Jacobian (actuator $\to$ task) \\
    $\bJ_h^\top$     & \cref{eq:Jh_T} & Maps EE wrench $\to$ actuator inputs \\
    $J_{\text{aug}}$ & \cref{eq:Jaug_explicit} & Augmented Jacobian (task + cable, $3\times 2$) \\
    $\dot{\bu}_v = \bJ_h^\dag\,\dot{\bp}$ & \cref{eq:uv_ik} & Hybrid velocity IK \\
    \bottomrule
  \end{tabular}
\end{center}

\noindent
Key relationships at a glance:
\begin{align*}
  \dot{\bp}  &= \bJ_h\,\dot{\bu}_v
               &&\text{(actuator velocity $\to$ EE velocity)} \\
  \bu        &= \bJ_h^\top \bF_{EE}
               &&\text{(EE force $\to$ actuator inputs, static)} \\
  \dot{l}    &= r_p\,\dot{q}_2
               &&\text{(cable rate from joint-2 rate)} \\
  \det\bJ_h  &= (L_1 L_2 / r_p)\sin q_2
               &&\text{(singular at } q_2 = 0, \pm\pi) \\
\end{align*}

% =============================================================================
\section{Cable Kinematics via Tangency Cancellation}
\label{sec:cable_kin_tangency}
% =============================================================================

This section derives the cable length rate $\dot{l}$ rigorously for a cable path
that wraps around the driven-pulley body before attaching to the load point.  The
key result is that the velocity of the intermediate tangency points cancels,
leaving a single clean expression involving only the load-attachment point.

\subsection{Assumptions and Setup}
\label{sec:cable_setup}

Assume the driven-pulley centre $O_2$ is \emph{fixed in the world}:
\begin{equation}
  {}^{W}\mathbf{r}_{WO_2} = \text{constant}
  \;\implies\;
  {}^{W}\mathbf{v}_{O_2} = {}^{W}\dot{\mathbf{r}}_{WO_2} = \mathbf{0}.
\end{equation}

The variable cable path is
\[
  B_{0,L} \;\to\; A_{1,L} \;\leadsto\; B_{1,L} \;\to\; A_{2,R},
\]
where:
\begin{itemize}
  \item $B_{0,L}$ --- fixed anchor point in the world frame.
  \item $A_{1,L}$ --- tangency point where the cable \emph{leaves} the straight
        segment and begins wrapping on the pulley.
  \item $B_{1,L}$ --- tangency point where the cable \emph{leaves} the pulley
        and begins the final straight segment.
  \item $A_{2,R}$ --- load-attachment point, rigidly fixed to the rotating body.
\end{itemize}

\subsection{Decomposition of Total Cable Length}
\label{sec:cable_decompose}

The total cable length is:
\begin{equation}
  l = l_{B_{0,L}A_{1,L}} + l_{A_{1,L}B_{1,L}} + l_{B_{1,L}A_{2,R}}.
\end{equation}

Differentiating:
\begin{equation}
  \dot{l}
  = \dot{l}_{B_{0,L}A_{1,L}}
  + \dot{l}_{A_{1,L}B_{1,L}}
  + \dot{l}_{B_{1,L}A_{2,R}}.
\end{equation}

\subsection{First Straight Segment}
\label{sec:seg1}

Define ${}^{W}\mathbf{s}_1 = {}^{W}\mathbf{r}_{WA_{1,L}} - {}^{W}\mathbf{r}_{WB_{0,L}}$.
Using $\frac{d}{dt}\|\mathbf{x}\| = \frac{\mathbf{x}^\top\dot{\mathbf{x}}}{\|\mathbf{x}\|}$
and noting that $B_{0,L}$ is fixed:
\begin{equation}
  \dot{l}_{B_{0,L}A_{1,L}}
  = {}^{W}\hat{\mathbf{u}}_{B_{0,L}A_{1,L}}^\top\,{}^{W}\dot{\mathbf{r}}_{WA_{1,L}},
  \qquad
  {}^{W}\hat{\mathbf{u}}_{B_{0,L}A_{1,L}}
  = \frac{{}^{W}\mathbf{r}_{WA_{1,L}} - {}^{W}\mathbf{r}_{WB_{0,L}}}
         {\left\|{}^{W}\mathbf{r}_{WA_{1,L}} - {}^{W}\mathbf{r}_{WB_{0,L}}\right\|}.
\end{equation}

\subsection{Wrapped Arc}
\label{sec:arc}

For the arc $A_{1,L} \to B_{1,L}$, define unit tangents ${}^{W}\mathbf{t}_{A_{1,L}}$
and ${}^{W}\mathbf{t}_{B_{1,L}}$ pointing \emph{outward from the arc} at each
endpoint (i.e.\ opposite to the direction of cable travel).  Then:
\begin{equation}
  \dot{l}_{A_{1,L}B_{1,L}}
  = +{}^{W}\mathbf{t}_{A_{1,L}}^\top\,{}^{W}\dot{\mathbf{r}}_{WA_{1,L}}
    - {}^{W}\mathbf{t}_{B_{1,L}}^\top\,{}^{W}\dot{\mathbf{r}}_{WB_{1,L}}.
\end{equation}

\noindent
\textit{Note on sign convention.} With the standard incoming-tangent convention
($\mathbf{t}$ pointing in the direction of travel), the signs would be reversed.
The outward convention used here ensures the tangency cancellation in
\cref{sec:tangency_cancel} follows with the same sign as the straight-segment
formula.

\subsection{Last Straight Segment}
\label{sec:seg2}

Define ${}^{W}\mathbf{s}_2 = {}^{W}\mathbf{r}_{WA_{2,R}} - {}^{W}\mathbf{r}_{WB_{1,L}}$:
\begin{equation}
  \dot{l}_{B_{1,L}A_{2,R}}
  = {}^{W}\hat{\mathbf{u}}_{B_{1,L}A_{2,R}}^\top
    \!\left({}^{W}\dot{\mathbf{r}}_{WA_{2,R}} - {}^{W}\dot{\mathbf{r}}_{WB_{1,L}}\right),
  \quad
  {}^{W}\hat{\mathbf{u}}_{B_{1,L}A_{2,R}}
  = \frac{{}^{W}\mathbf{r}_{WA_{2,R}} - {}^{W}\mathbf{r}_{WB_{1,L}}}
         {\left\|{}^{W}\mathbf{r}_{WA_{2,R}} - {}^{W}\mathbf{r}_{WB_{1,L}}\right\|}.
\end{equation}

\subsection{Summing All Three Contributions}
\label{sec:sum}

Adding the three parts:
\begin{align}
  \dot{l}
  &=
  \underbrace{\left(
    {}^{W}\hat{\mathbf{u}}_{B_{0,L}A_{1,L}}
    + {}^{W}\mathbf{t}_{A_{1,L}}
  \right)^\top}_{\text{coefficient of }\dot{\mathbf{r}}_{WA_{1,L}}}
  {}^{W}\dot{\mathbf{r}}_{WA_{1,L}}
  +\;
  {}^{W}\hat{\mathbf{u}}_{B_{1,L}A_{2,R}}^\top\,{}^{W}\dot{\mathbf{r}}_{WA_{2,R}}
  \nonumber\\[4pt]
  &\quad
  -\underbrace{\left(
    {}^{W}\mathbf{t}_{B_{1,L}}
    + {}^{W}\hat{\mathbf{u}}_{B_{1,L}A_{2,R}}
  \right)^\top}_{\text{coefficient of }{}^{W}\dot{\mathbf{r}}_{WB_{1,L}}}
  {}^{W}\dot{\mathbf{r}}_{WB_{1,L}}.
\end{align}

\subsection{Tangency Cancellation}
\label{sec:tangency_cancel}

At each tangency point the cable direction (straight segment unit vector) must
equal the outward arc tangent for the cable to be smooth (no kink).  Therefore:
\begin{equation}
  {}^{W}\mathbf{t}_{A_{1,L}} = -{}^{W}\hat{\mathbf{u}}_{B_{0,L}A_{1,L}},
  \qquad
  {}^{W}\mathbf{t}_{B_{1,L}} = -{}^{W}\hat{\mathbf{u}}_{B_{1,L}A_{2,R}}.
\end{equation}

Both bracket terms vanish identically:
\begin{equation}
  {}^{W}\hat{\mathbf{u}}_{B_{0,L}A_{1,L}} + {}^{W}\mathbf{t}_{A_{1,L}} = \mathbf{0},
  \qquad
  {}^{W}\mathbf{t}_{B_{1,L}} + {}^{W}\hat{\mathbf{u}}_{B_{1,L}A_{2,R}} = \mathbf{0}.
\end{equation}

The velocities of the two tangency points $A_{1,L}$ and $B_{1,L}$ drop out completely,
yielding:
\begin{equation}
  \label{eq:ldot_key}
  \boxed{
    \dot{l}
    = {}^{W}\hat{\mathbf{u}}_{B_{1,L}A_{2,R}}^\top\;{}^{W}\dot{\mathbf{r}}_{WA_{2,R}}
  }
\end{equation}

\subsection{Rigid-Body Velocity of the Load-Attachment Point}
\label{sec:rigid_body}

Because $A_{2,R}$ is fixed to the body rotating about the fixed point $O_2$, and
${}^{W}\dot{\mathbf{r}}_{WO_2} = \mathbf{0}$:
\begin{equation}
  {}^{W}\dot{\mathbf{r}}_{WA_{2,R}}
  = {}^{W}\boldsymbol{\omega}_2 \times {}^{W}\mathbf{r}_{O_2 A_{2,R}}.
\end{equation}

Substituting into \cref{eq:ldot_key}:
\begin{equation}
  \label{eq:ldot_omega}
  \boxed{
    \dot{l}
    = {}^{W}\hat{\mathbf{u}}_{B_{1,L}A_{2,R}}^\top
      \left(
        {}^{W}\boldsymbol{\omega}_2
        \times
        {}^{W}\mathbf{r}_{O_2 A_{2,R}}
      \right)
  }
\end{equation}

\subsection{Planar Reduction}
\label{sec:planar}

For planar motion, ${}^{W}\boldsymbol{\omega}_2 = [0,\;0,\;\dot{\theta}_2]^\top$,
giving:
\begin{equation}
  \label{eq:ldot_planar}
  \boxed{
    \dot{l}
    = {}^{W}\hat{\mathbf{u}}_{B_{1,L}A_{2,R}}^\top
      \left(
        \begin{bmatrix} 0 \\ 0 \\ \dot{\theta}_2 \end{bmatrix}
        \times
        {}^{W}\mathbf{r}_{O_2 A_{2,R}}
      \right)
    \;=\;
    J_l(\theta_2)\,\dot{\theta}_2 .
  }
\end{equation}

The scalar Jacobian entry $J_l(\theta_2)$ depends only on the current geometry of
the cable's last straight segment and the lever arm ${}^{W}\mathbf{r}_{O_2 A_{2,R}}$;
it replaces the constant $r_p$ approximation used in the simplified model of
\cref{sec:act_j2}.  When the cable wrapping is symmetric and $A_{2,R}$ lies on
the pitch circle, $J_l = r_p$, recovering the belt formula \cref{eq:cable_kin}.

\subsection{Spatial-Velocity Row Jacobian Form}
\label{sec:spatial_jac}

The result \cref{eq:ldot_omega} can be rewritten in the compact \emph{spatial velocity}
form that appears in cable-robot literature, making the translational and rotational
contributions explicit side-by-side.

\paragraph{Triple-product identity.}
For any vectors $\mathbf{u}$, $\boldsymbol{\omega}$, $\mathbf{r}$:
\begin{equation}
  \label{eq:triple}
  \mathbf{u}^\top(\boldsymbol{\omega} \times \mathbf{r})
  = (\mathbf{r} \times \mathbf{u})^\top \boldsymbol{\omega}.
\end{equation}

Apply this to \cref{eq:ldot_omega} with
$\mathbf{u} \triangleq {}^{W}\hat{\mathbf{u}}_{B_{1,L}A_{2,R}}$ and
$\mathbf{r} \triangleq {}^{W}\mathbf{r}_{O_2 A_{2,R}}$:

\begin{equation}
  \dot{l}
  = \mathbf{u}^\top \underbrace{{}^{W}\mathbf{v}_{O_2}}_{\mathbf{0}}
  + \mathbf{u}^\top({}^{W}\boldsymbol{\omega}_2 \times \mathbf{r})
  = \mathbf{u}^\top\,{}^{W}\mathbf{v}_{O_2}
  + (\mathbf{r} \times \mathbf{u})^\top\,{}^{W}\boldsymbol{\omega}_2 .
\end{equation}

Writing this as a single row times a spatial velocity vector:

\begin{equation}
  \label{eq:spatial_row}
  \boxed{
  \dot{l}
  =
  \underbrace{%
    \begin{bmatrix}
      {}^{W}\hat{\mathbf{u}}_{B_{1,L}A_{2,R}}^\top &
      \bigl({}^{W}\mathbf{r}_{O_2 A_{2,R}}
            \times {}^{W}\hat{\mathbf{u}}_{B_{1,L}A_{2,R}}\bigr)^\top
    \end{bmatrix}
  }_{V \;\in\;\mathbb{R}^{1\times 6}}
  \begin{bmatrix}
    {}^{W}\mathbf{v}_{O_2} \\[3pt] {}^{W}\boldsymbol{\omega}_2
  \end{bmatrix}
  }
\end{equation}

The first block ${}^{W}\hat{\mathbf{u}}^\top \in \mathbb{R}^{1\times 3}$ captures the
\emph{translational} contribution; the second block
$(\mathbf{r}\times\mathbf{u})^\top \in \mathbb{R}^{1\times 3}$ captures the
\emph{rotational} (moment-arm) contribution.  This is structurally identical to the
structure matrix row in cable-driven parallel robot formulations.

\paragraph{Fixed-pivot reduction.}
Since $O_2$ is fixed, ${}^{W}\mathbf{v}_{O_2} = \mathbf{0}$, and \cref{eq:spatial_row}
reduces to:

\begin{equation}
  \label{eq:rot_only}
  \boxed{
  \dot{l}
  =
  \bigl({}^{W}\mathbf{r}_{O_2 A_{2,R}}
        \times {}^{W}\hat{\mathbf{u}}_{B_{1,L}A_{2,R}}\bigr)^\top
  {}^{W}\boldsymbol{\omega}_2
  }
\end{equation}

\noindent
This is the rotational-only form: the cable rate equals the moment-arm vector
(lever arm crossed with cable unit vector) dotted with the angular velocity of the
driven link.

\subsection{Multi-Cable Structure Matrix}
\label{sec:structure_matrix}

For a system with $m$ cables, each cable $i$ contributes one row.  Stack all rows
to obtain the \emph{structure matrix} $V \in \mathbb{R}^{m \times 6}$:

\begin{equation}
  \label{eq:structure_matrix}
  \boxed{
  V
  =
  \begin{bmatrix}
    {}^{W}\hat{\mathbf{u}}_1^\top &
    \bigl({}^{W}\mathbf{r}_1 \times {}^{W}\hat{\mathbf{u}}_1\bigr)^\top \\[4pt]
    {}^{W}\hat{\mathbf{u}}_2^\top &
    \bigl({}^{W}\mathbf{r}_2 \times {}^{W}\hat{\mathbf{u}}_2\bigr)^\top \\[2pt]
    \vdots & \vdots \\[2pt]
    {}^{W}\hat{\mathbf{u}}_m^\top &
    \bigl({}^{W}\mathbf{r}_m \times {}^{W}\hat{\mathbf{u}}_m\bigr)^\top
  \end{bmatrix},
  \qquad
  \dot{\mathbf{l}}
  = V
  \begin{bmatrix}
    {}^{W}\mathbf{v} \\ {}^{W}\boldsymbol{\omega}
  \end{bmatrix},
  }
\end{equation}

\noindent
where ${}^{W}\hat{\mathbf{u}}_i \triangleq {}^{W}\hat{\mathbf{u}}_{B_{1,L}^{(i)}A_{2,R}^{(i)}}$
and ${}^{W}\mathbf{r}_i \triangleq {}^{W}\mathbf{r}_{O_2 A_{2,R}^{(i)}}$ for the $i$-th
cable, and the spatial velocity $[{}^{W}\mathbf{v}^\top,\;{}^{W}\boldsymbol{\omega}^\top]^\top$
is evaluated at the pivot point $O_2$.

\paragraph{Connection to the hybrid Jacobian.}
For the single cable driving joint~2 of the planar arm the angular velocity is
${}^{W}\boldsymbol{\omega}_2 = \dot{q}_2\,\hat{\mathbf{k}}$ with
$\hat{\mathbf{k}} = [0,0,1]^\top$.
Substituting into \cref{eq:rot_only}:

\begin{equation}
  \dot{l}
  = \bigl({}^{W}\mathbf{r}_{O_2 A_{2,R}}
          \times {}^{W}\hat{\mathbf{u}}_{B_{1,L}A_{2,R}}\bigr)^\top
    \hat{\mathbf{k}}\;\dot{q}_2
  = J_l(q_2)\,\dot{q}_2 ,
\end{equation}

recovering precisely the scalar Jacobian entry from \cref{eq:ldot_planar}.

\paragraph{Derivation of the belt formula $J_l = r_p$.}
We now show explicitly that the tangency geometry forces $J_l = r_p$.

\medskip\noindent
\textit{Setup.}
Let $\phi$ be the angle that the lever arm makes with the world $x$-axis at an
arbitrary cable configuration.  With $A_{2,R}$ on the pitch circle of radius $r_p$:
\begin{equation*}
  {}^{W}\mathbf{r}_{O_2 A_{2,R}}
  = r_p
    \begin{bmatrix} \cos\phi \\ \sin\phi \\ 0 \end{bmatrix}.
\end{equation*}

\textit{Tangency condition.}
Because the cable departs the pulley tangentially, the cable direction must be
perpendicular to the lever arm:
${}^{W}\hat{\mathbf{u}} \perp {}^{W}\mathbf{r}_{O_2 A_{2,R}}$, i.e.\ the angle
between them is $90°$.  The unique planar unit tangent is:
\begin{equation*}
  {}^{W}\hat{\mathbf{u}}_{B_{1,L}A_{2,R}}
  = \begin{bmatrix} -\sin\phi \\ \cos\phi \\ 0 \end{bmatrix}
  \qquad\text{(or its negation, depending on winding direction)}.
\end{equation*}

\textit{Cross product.}
\begin{equation}
  \label{eq:cross_belt}
  {}^{W}\mathbf{r}_{O_2 A_{2,R}} \times {}^{W}\hat{\mathbf{u}}_{B_{1,L}A_{2,R}}
  =
  r_p
  \begin{bmatrix} \cos\phi \\ \sin\phi \\ 0 \end{bmatrix}
  \times
  \begin{bmatrix} -\sin\phi \\ \cos\phi \\ 0 \end{bmatrix}
  =
  r_p
  \begin{bmatrix} 0 \\ 0 \\ \cos^2\!\phi + \sin^2\!\phi \end{bmatrix}
  =
  r_p\,\hat{\mathbf{k}}.
\end{equation}

The $\hat{\mathbf{k}}$ component collapses to $1$ by the Pythagorean identity regardless
of $\phi$.

\textit{Projecting onto $\hat{\mathbf{k}}$.}
\begin{equation}
  \label{eq:belt_derivation}
  J_l
  = \bigl({}^{W}\mathbf{r}_{O_2 A_{2,R}} \times {}^{W}\hat{\mathbf{u}}\bigr)^\top\hat{\mathbf{k}}
  = \bigl(r_p\,\hat{\mathbf{k}}\bigr)^\top\hat{\mathbf{k}}
  = r_p\,\underbrace{\hat{\mathbf{k}}^\top\hat{\mathbf{k}}}_{=\,1}
  = r_p .
\end{equation}

\begin{equation}
  \label{eq:belt_final}
  \boxed{J_l = r_p}
\end{equation}

This is the belt (pitch-circle) formula: the effective moment arm of a cable
tangent to a pulley of pitch radius $r_p$ equals $r_p$ for every cable angle
$\phi$, which is exactly the constant mapping encoded in the actuation matrix
$\bA = \operatorname{diag}(1, r_p)$ from \cref{eq:act_map}.

\medskip\noindent
\textit{Geometric intuition.}
The cross-product magnitude $|\mathbf{r}\times\hat{\mathbf{u}}| = |\mathbf{r}||\hat{\mathbf{u}}|\sin\alpha$,
where $\alpha$ is the angle between the lever arm and cable direction.
Tangency forces $\alpha = 90°$; hence $\sin\alpha = 1$ and the magnitude locks
to $|\mathbf{r}| = r_p$ regardless of joint angle $q_2$.  The spatial-velocity
framework makes this structure manifest: it is built into the cross-product term
of every row of the structure matrix $V$.

\begin{mdframed}
\begin{corollary}[Attachment-point invariance]
\label{cor:attachment_invariance}
Moving the cable attachment point $A_{2,R}$ to any other location on the
same pitch circle — i.e.\ changing the angle $\phi$ while keeping
$|{}^W\mathbf{r}_{O_2 A_{2,R}}| = r_p$ — does \textbf{not} change $J_l$.
Equivalently, $J_l = r_p$ is independent of
(i) the joint angle $q_2$, and
(ii) the circumferential position of $A_{2,R}$ on the pulley.
\end{corollary}
\end{mdframed}

\begin{proof}
By \cref{eq:cross_belt}, the cross product evaluates to
$r_p(\cos^2\phi + \sin^2\phi)\,\hat{\mathbf{k}} = r_p\,\hat{\mathbf{k}}$.
The Pythagorean identity eliminates $\phi$ completely, so the projection
onto $\hat{\mathbf{k}}$ in \cref{eq:belt_derivation} gives $J_l = r_p$
for all $\phi$.
\end{proof}

\noindent
\textbf{Consequence for re-design.}
The three cases below follow directly from the corollary:
\begin{itemize}[leftmargin=2em, topsep=2pt, itemsep=2pt]
  \item \emph{Different $\phi$, same $r_p$}: rotating or repositioning $A_{2,R}$ around the
        pulley centre leaves $J_l$ unchanged.  Only the cable path geometry
        (tangent angle, wrap arc) changes.
  \item \emph{Different $q_2$}: as the link rotates, $A_{2,R}$ moves in world
        frame but remains on the pitch circle.  $J_l$ stays constant — this
        is why the actuation matrix $\bA = \operatorname{diag}(1,r_p)$ is
        pose-independent.
  \item \emph{Different $r_p$}: swapping to a pulley of radius $r_p'$ scales
        $J_l \to r_p'$ and proportionally rescales the entire hybrid-Jacobian
        column $\bJ_h^{(2)} = \bJ^{(2)}/r_p$.
\end{itemize}

% =============================================================================
\section{Concrete Velocity Mappings for the CupManipulatorTendon}
\label{sec:vel_mappings}
% =============================================================================

This section collects, in one place, all three velocity-level mappings that are
used in the Python implementation:
\begin{enumerate}
  \item $\dot{\bq} \;\to\; \dot{l}$ \quad (cable length rate, single cable to joint 2)
  \item $\dot{\bq} \;\to\; \dot{\bp}$ \quad (end-effector Cartesian velocity)
  \item $\dot{l},\;\dot{q}_1 \;\to\; \dot{\bp}$ \quad (actuation-space to EE velocity)
\end{enumerate}

Numerical values used throughout:
\[
  L_1 = 0.34247\;\text{m}, \quad
  L_2 = 0.19000\;\text{m}, \quad
  r_p = 0.04775\;\text{m}.
\]

% ---------------------------------------------------------------------------
\subsection{Joint-Velocity to Cable-Length Rates}
\label{sec:qdot2ldot}
% ---------------------------------------------------------------------------

Joint~2 is driven by \textbf{two cables}, green ($G$) and red ($R$), which wrap
on the big pulley from \emph{opposite sides} (see \cref{sec:antagonistic}).
A positive $\dot{q}_2$ simultaneously \emph{pays out} the green cable and
\emph{takes up} the red cable.  From the belt formula \cref{eq:belt_final}
applied to each cable:

\begin{equation}
  \label{eq:ldot_q}
  \boxed{
    \dot{l}_G = +r_p\,\dot{q}_2,
    \qquad
    \dot{l}_R = -r_p\,\dot{q}_2 .
  }
\end{equation}

The two rates are not independent: they satisfy the
\emph{antagonistic constraint} $\dot{l}_G + \dot{l}_R = 0$ at all times
(\cref{eq:antag_constraint}).  Therefore only \textbf{one scalar}
$\dot{l} \triangleq \dot{l}_G = -\dot{l}_R$ is needed to fully describe the
cable kinematics of joint~2.

\noindent
\textbf{Dependencies:} $\dot{l}_G$ and $\dot{l}_R$ depend \emph{only} on
$\dot{q}_2$; they are independent of $q_1$, $q_2$, and $\dot{q}_1$.

In vector form for the full joint-velocity vector
$\dot{\bq} = [\dot{q}_1,\,\dot{q}_2]^\top$:

\begin{equation}
  \label{eq:ldot_row}
  \bl{\dot{l}_G \choose \dot{l}_R}
  =
  \underbrace{r_p\begin{bmatrix} 0 & +1 \\ 0 & -1 \end{bmatrix}}_{\bm{J}_l \;\in\;\mathbb{R}^{2\times 2}}
  \begin{bmatrix} \dot{q}_1 \\ \dot{q}_2 \end{bmatrix}.
\end{equation}

Because the rows are linearly dependent (rank~1), the equivalent
\emph{scalar} row Jacobian using only $\dot{l} \triangleq \dot{l}_G$ is:
\begin{equation}
  \label{eq:ldot_scalar}
  \dot{l}_G
  = \underbrace{\begin{bmatrix} 0 & r_p \end{bmatrix}}_{\bm{J}_{l,G} \;\in\;\mathbb{R}^{1\times 2}}
  \begin{bmatrix} \dot{q}_1 \\ \dot{q}_2 \end{bmatrix}
  = 0.04775\;\dot{q}_2 \quad [\text{m/s per rad/s}].
\end{equation}

\noindent
The hybrid Jacobian (\cref{sec:actspc2ee}) uses $\dot{l} \equiv \dot{l}_G$ as
the second actuation-space coordinate; $\dot{l}_R$ is then implicit via
$\dot{l}_R = -\dot{l}_G$.

% ---------------------------------------------------------------------------
\subsection{Joint-Velocity to End-Effector Velocity}
\label{sec:qdot2ee}
% ---------------------------------------------------------------------------

From the geometric Jacobian \cref{eq:Jgeo}:

\begin{equation}
  \label{eq:ee_vel}
  \boxed{
  \begin{bmatrix} \dot{p}_x \\ \dot{p}_y \end{bmatrix}
  =
  \begin{bmatrix}
    -L_1 s_1 - L_2 s_{12} & -L_2 s_{12} \\[4pt]
     L_1 c_1 + L_2 c_{12} &  L_2 c_{12}
  \end{bmatrix}
  \begin{bmatrix} \dot{q}_1 \\ \dot{q}_2 \end{bmatrix}
  }
\end{equation}

\noindent
Expanding the two rows explicitly:

\begin{align}
  \dot{p}_x &= -(L_1 s_1 + L_2 s_{12})\,\dot{q}_1 \;-\; L_2 s_{12}\,\dot{q}_2 ,
  \label{eq:px_dot}\\
  \dot{p}_y &= \phantom{-}(L_1 c_1 + L_2 c_{12})\,\dot{q}_1 \;+\; L_2 c_{12}\,\dot{q}_2 .
  \label{eq:py_dot}
\end{align}

\paragraph{Observation.}
$\dot{q}_1$ controls both $\dot{p}_x$ and $\dot{p}_y$ via the full arm reach
$(L_1 + L_2)$, whereas $\dot{q}_2$ controls them only via the distal link $L_2$.
Near $q_2 = 0$ (arm straight), the second column of $\bJ$ becomes parallel to
the first, causing $\det\bJ = L_1 L_2 \sin q_2 \to 0$.

% ---------------------------------------------------------------------------
\subsection{Actuation-Space to EE Velocity (Hybrid Form)}
\label{sec:actspc2ee}
% ---------------------------------------------------------------------------

The hybrid Jacobian from \cref{eq:Jh} maps the \emph{actuation-space velocity}
$\dot{\bu}_v = [\dot{q}_1,\,\dot{l}_G]^\top$ to EE velocity
(recall $\dot{l}_G = -\dot{l}_R$ so $\dot{l}_R$ adds no new information):

\begin{equation}
  \label{eq:hybrid_vel_exp}
  \dot{\bp} = \bJ_h\,\dot{\bu}_v
  \quad\Longleftrightarrow\quad
  \begin{bmatrix} \dot{p}_x \\ \dot{p}_y \end{bmatrix}
  =
  \underbrace{%
  \begin{bmatrix}
    -L_1 s_1 - L_2 s_{12} & -\dfrac{L_2 s_{12}}{r_p} \\[8pt]
     L_1 c_1 + L_2 c_{12} &  \dfrac{L_2 c_{12}}{r_p}
  \end{bmatrix}
  }_{\bJ_h(\bq)}
  \begin{bmatrix} \dot{q}_1 \\ \dot{l}_G \end{bmatrix}.
\end{equation}

Numerically, $L_2 / r_p = 0.190 / 0.04775 \approx 3.979$, so the cable column
amplifies EE velocity by a factor of $\approx 4$ compared to a direct-drive
joint of the same link length.

% ---------------------------------------------------------------------------
\subsection{Complete Mapping Summary}
\label{sec:mapping_summary}
% ---------------------------------------------------------------------------

\begin{center}
\renewcommand{\arraystretch}{1.6}
\begin{tabular}{p{5.1cm}ll}
  \toprule
  Mapping & Equation & Notes \\
  \midrule
  $\dot{q}_2 \to \dot{l}_G,\,\dot{l}_R$
    & $\dot{l}_G = +r_p\dot{q}_2,\;\dot{l}_R = -r_p\dot{q}_2$
    & Antagonistic pair; $\dot{l}_G + \dot{l}_R = 0$ \\
  $\dot{q}_2 \to \dot{l}_G$ (scalar)
    & $\dot{l}_G = r_p\,\dot{q}_2$
    & Used as the independent cable rate \\
  $\dot{\bq} \to \dot{\bp}$ (joint space)
    & $\dot{\bp} = \bJ(\bq)\,\dot{\bq}$
    & Standard geometric Jacobian \cref{eq:Jgeo} \\
  $\dot{\bq} \to \dot{\bp}$ (actuation space)
    & $\dot{\bp} = \bJ_h(\bq)\,\dot{\bu}_v$
    & $\dot{\bu}_v = [\dot{q}_1,\,\dot{l}_G]^\top$; $\dot{l}_R$ implicit \\
  $\dot{\bp} \to \dot{\bq}$ (velocity IK)
    & $\dot{\bq} = \bJ^{\dag}\dot{\bp}$
    & Damped pseudo-inverse, singular near $q_2=0$ \\
  $\dot{\bp} \to \dot{\bu}_v$ (hybrid vel IK)
    & $\dot{\bu}_v = \bJ_h^{\dag}\dot{\bp}$
    & Used in \texttt{compute\_velocity\_ik()} \\
  \bottomrule
\end{tabular}
\end{center}

\paragraph{Python implementation.}
In \texttt{robots/cup\_manipulator\_tendon.py}:
\begin{itemize}[leftmargin=2em]
  \item \texttt{compute\_ik\_analytical(plant, target\_xy, ...)} solves
        $\dot{\bq} \to \bp$ by the closed-form 2R formula (preferred).
  \item \texttt{compute\_velocity\_ik(plant, context, ee\_velocity\_xy, damping)}
        computes $\dot{\bu}_v = \bJ_h^\dag\,\dot{\bp}$ via the damped pseudo-inverse.
        $\dot{l}_G$ is then recovered as $\dot{l}_G = (\dot{\bu}_v)_2$,
        $\dot{q}_2 = \dot{l}_G/r_p$, and
        $\dot{l}_R = -\dot{l}_G$ (antagonistic constraint).
\end{itemize}

% ---------------------------------------------------------------------------
\subsection{Antagonistic Cable Pair}
\label{sec:antagonistic}
% ---------------------------------------------------------------------------

Joint~2 is driven by \textbf{two cables that act in opposition} --- analogous to
agonist/antagonist muscle pairs in biology.  Each cable can only \emph{pull}
(tension $\ge 0$); the antagonistic arrangement allows bidirectional torque on a
passive revolute joint.

\paragraph{Physical routing (from \texttt{CableRig}).}
\begin{center}
\renewcommand{\arraystretch}{1.3}
\begin{tabular}{lll}
  \toprule
  Cable & Colour & Route through pulleys \\
  \midrule
  Agonist   & \textbf{Green} &
    Start$_R$ $\to$ Drive$_{B_R}$ $\xrightarrow{\text{wrap}}$ IdlerR
    $\xrightarrow{\text{wrap}}$ BigPulley$_{A_L \to B_L}$ $\to$ EndL \\
  Antagonist & \textbf{Red} &
    Start$_L$ $\to$ Drive$_{B_L}$ $\xrightarrow{\text{wrap}}$ IdlerL
    $\xrightarrow{\text{wrap}}$ BigPulley$_{A_R \to B_R}$ $\to$ EndR \\
  \bottomrule
\end{tabular}
\end{center}

The two cables enter the \emph{big pulley} (rigidly attached to link~2,
body \texttt{link2\_tendon}) from \textbf{opposite sides}: green from the left
($A_L$) and red from the right ($A_R$).  This means they wrap in opposing angular
senses, producing opposite torques on joint~2.

\paragraph{Cable-rate kinematics of the antagonistic pair.}
Let $l_G$ and $l_R$ denote the \emph{free lengths} (drive-to-anchor lengths) of
the green and red cables respectively.  Since the wrap arcs on the big pulley are
complementary, a rotation $\dot{q}_2$ simultaneously pays out one cable and takes
up the other:

\begin{equation}
  \label{eq:ldot_pair}
  \dot{l}_G = +r_p\,\dot{q}_2,
  \qquad
  \dot{l}_R = -r_p\,\dot{q}_2.
\end{equation}

Summing the two rates gives the \emph{antagonistic constraint}:

\begin{equation}
  \label{eq:antag_constraint}
  \boxed{\dot{l}_G + \dot{l}_R = 0}
\end{equation}

This is the cable equivalent of muscle incompressibility: whatever length the
agonist gains, the antagonist gives up.

In matrix form for $\dot{\bl} = [\dot{l}_G,\,\dot{l}_R]^\top$:

\begin{equation}
  \label{eq:Jl_pair}
  \dot{\bl}
  =
  \underbrace{%
    r_p
    \begin{bmatrix} 0 & +1 \\ 0 & -1 \end{bmatrix}
  }_{\bm{J}_l^{(2)} \;\in\;\mathbb{R}^{2\times 2}}
  \begin{bmatrix} \dot{q}_1 \\ \dot{q}_2 \end{bmatrix}.
\end{equation}

\paragraph{Net torque on joint 2.}
Each cable exerts a tension force on the big pulley with a moment arm of $r_p$.
Because the cables wrap in opposite senses, green tension $f_G \ge 0$ produces
positive torque and red tension $f_R \ge 0$ produces negative torque (or vice
versa, depending on the positive-$q_2$ convention):

\begin{equation}
  \label{eq:tau2_pair}
  \boxed{
    \tau_2 = r_p\,(f_G - f_R),
    \qquad f_G,\,f_R \ge 0 .
  }
\end{equation}

This matches the structure-matrix transposition: the wrench on joint~2 from the
cable-force vector $\bf_c = [f_G,\,f_R]^\top$ is $\btau = (\bm{J}_l^{(2)})^\top\!\bf_c$,
giving $\tau_2 = r_p(f_G - f_R)$ and $\tau_1 = 0$ (cables do not act on joint~1).

\paragraph{Drive-motor coupling.}
The drive pulley (on the base, body \texttt{link1\_base}) rotates by
$\theta_{\mathrm{drv}}$.  Its two exits feed the green and red cables with
\emph{equal and opposite} length changes (the drive pulley is also symmetric):

\begin{equation}
  \label{eq:ldot_drive}
  \dot{l}_G^{\mathrm{drv}} = -r_{\mathrm{drv}}\,\dot{\theta}_{\mathrm{drv}},
  \qquad
  \dot{l}_R^{\mathrm{drv}} = +r_{\mathrm{drv}}\,\dot{\theta}_{\mathrm{drv}},
\end{equation}

where $r_{\mathrm{drv}} = r_p$ (same HTD 5M tooth pitch).  The net free length
of each cable is determined by the \emph{difference} of big-pulley and
drive-pulley contributions.  In the implementation, the drive motor is the
sole control input $u_2$ for joint~2; the antagonistic constraint ensures that
setting $\dot{\theta}_{\mathrm{drv}}$ is sufficient to specify both $\dot{l}_G$
and $\dot{l}_R$ simultaneously.

\paragraph{Admissible-force cone.}
Since both cable tensions must be non-negative, the net torque is bounded:

\begin{equation}
  \label{eq:torque_cone}
  \tau_2 \in [-r_p f_{\max},\; +r_p f_{\max}],
\end{equation}

but the pair can also generate a \emph{controlled stiffness}: choosing
$f_G = f_R = f_{\text{co}} > 0$ (co-contraction) produces $\tau_2 = 0$ while
increasing the effective joint stiffness.  This is the cable analogue of
muscle co-activation.

% =============================================================================
\begin{thebibliography}{9}
\bibitem{pott2018cable}
  A.~Pott and T.~Bruckmann,
  \textit{Cable-Driven Parallel Robots: Theory and Application},
  Springer Tracts in Advanced Robotics, 2018.
\end{thebibliography}

\end{document}
